% GENERATED FILE. Edit canonical lessons, then run npm run generate:books.
% canonical-source-hash: fnv1a64:46be862cd092f590
% canonical-lessons: ZH-C495-qizhong, ZH-C495-kaochang, ZH-C495-manfen, ZH-C495-zuobi, ZH-C495-nanti

\chapter{Marks}
\label{ch:zh-marks}
\hlchaptermodality{495}

\begin{chapteropening}
\textbf{By the end of this chapter:} \emph{I can talk about a midterm exam, the exam hall, full marks, cheating and hard questions.}

Complete the last lesson of chapter 495: I can talk about a midterm exam, the exam hall, full marks, cheating and hard questions.
\end{chapteropening}

\section[\texorpdfstring{\zh{期中} (qīzhōng)}{期中 (qīzhōng)}]{\zh{期中} (qīzhōng) — the middle of the term; midterm}

\label{lesson:ZH-C495-qizhong}



\begin{quote}
Before the new one: say the Mandarin for to pass the entrance exam for, then the Mandarin for to fail an entrance exam.
\end{quote}

\subsection*{You'll want to know: \zh{期中}}
\textbf{\zh{期中}} — \emph{qīzhōng} — \textquotedblleft{}the middle of the term; midterm\textquotedblright{}.

\begin{quote}
\emph{zhè cì qīzhōng kǎoshì wǒ gǎnjué hái xíng} — I think this midterm went all right.
\end{quote}

\begin{cousinweb}
\textbf{\zh{期中}} (\emph{qīzhōng}) is built from \zh{期} 'a set time' and \zh{中} 'the middle', both inherited from Middle Chinese, \emph{gi} and \emph{trjuwng} (Guangyun, in Baxter's transcription), which glosses the first 'an appointed time'. The modern first tone of qī is irregular: \emph{gi} would regularly give qí.
\end{cousinweb}

\begin{grammarlens}[title={What you've built}]
One more word for saying how it was.
\end{grammarlens}

\subsection*{Your turn}
\begin{itemize}
\raggedright
  \item \emph{Say it:} \emph{qīzhōng}
  \item \emph{Say it:} \emph{qīzhōng}, once more
  \item \emph{Say it:} \emph{zhè cì qīzhōng kǎoshì wǒ gǎnjué hái xíng} — I think this midterm went all right.
  \item \emph{Recall:} say the Mandarin for to pass the entrance exam for, then the Mandarin for to fail an entrance exam, then say \emph{qīzhōng} again
\end{itemize}

\begin{tcolorbox}[breakable,colback=teal!4,colframe=teal!35!black,title={Before you move on}]
What does \zh{期中} mean? (\textquotedblleft{}The middle of the term; midterm\textquotedblright{}.) Say it once more.
\end{tcolorbox}
\section[\texorpdfstring{\zh{考场} (kǎochǎng)}{考场 (kǎochǎng)}]{\zh{考场} (kǎochǎng) — exam hall; exam room}

\label{lesson:ZH-C495-kaochang}



\begin{quote}
Before the new one: say the Mandarin for to fail an entrance exam, then the Mandarin for the middle of the term.
\end{quote}

\subsection*{You'll want to know: \zh{考场}}
\textbf{\zh{考场}} — \emph{kǎochǎng} — \textquotedblleft{}exam hall; exam room\textquotedblright{}.

\begin{quote}
\textbf{\zh{进考场以前我喝了一杯水}} \emph{jìn kǎochǎng yǐqián wǒ hē le yì bēi shuǐ} — Before going into the exam hall, I drank a glass of water.
\end{quote}

Two third tones meet in this word: kǎo is said with a rising tone before chǎng, and the spelling keeps both marks.

\begin{cousinweb}
\textbf{\zh{考场}} (\emph{kǎochǎng}) is built from \zh{考} 'to examine' and \zh{场} 'a ground', both inherited from Middle Chinese, \emph{khawX} and \emph{drjang} (Guangyun, in Baxter's transcription), which glosses the second 'a threshing floor'; the third tone of chǎng is irregular for \emph{drjang}.
\end{cousinweb}

\begin{grammarlens}[title={What you've built}]
One more word for saying how it was.
\end{grammarlens}

\subsection*{Your turn}
\begin{itemize}
\raggedright
  \item \emph{Say it:} \emph{kǎochǎng}
  \item \emph{Say it:} \emph{kǎochǎng}, once more
  \item \emph{Say it:} \emph{jìn kǎochǎng yǐqián wǒ hē le yì bēi shuǐ} — Before going into the exam hall, I drank a glass of water.
  \item \emph{Recall:} say the Mandarin for to fail an entrance exam, then the Mandarin for the middle of the term, then say \emph{kǎochǎng} again
\end{itemize}

\begin{tcolorbox}[breakable,colback=teal!4,colframe=teal!35!black,title={Before you move on}]
What does \zh{考场} mean? (\textquotedblleft{}Exam hall; exam room\textquotedblright{}.) Say it once more.
\end{tcolorbox}
\section[mǎnfēn]{mǎnfēn — full marks}

\label{lesson:ZH-C495-manfen}



\begin{quote}
Before the new one: say the Mandarin for the middle of the term, then the Mandarin for an exam hall.
\end{quote}

\subsection*{You'll want to know: mǎnfēn}
\emph{mǎnfēn} — \textquotedblleft{}full marks\textquotedblright{}.

\begin{quote}
\emph{zhè cì shùxué kǎoshì tā shì mǎnfēn} — He got full marks in this maths test.
\end{quote}

Its characters come later; for now it is a word for the ear.

\begin{cousinweb}
\emph{mǎnfēn} is built from \emph{mǎn} 'full' and \emph{fēn} 'a part', both inherited from Middle Chinese, \emph{manX} and \emph{pjun} (Guangyun, in Baxter's transcription), which glosses the first 'brimming' and the second 'to divide'.
\end{cousinweb}

\begin{grammarlens}[title={What you've built}]
One more word for saying how it was.
\end{grammarlens}

\subsection*{Your turn}
\begin{itemize}
\raggedright
  \item \emph{Say it:} \emph{mǎnfēn}
  \item \emph{Say it:} \emph{mǎnfēn}, once more
  \item \emph{Say it:} \emph{zhè cì shùxué kǎoshì tā shì mǎnfēn} — He got full marks in this maths test.
  \item \emph{Recall:} say the Mandarin for the middle of the term, then the Mandarin for an exam hall, then say \emph{mǎnfēn} again
\end{itemize}

\begin{tcolorbox}[breakable,colback=teal!4,colframe=teal!35!black,title={Before you move on}]
What does mǎnfēn mean? (\textquotedblleft{}Full marks\textquotedblright{}.) Say it once more.
\end{tcolorbox}
\section[zuòbì]{zuòbì — to cheat in an exam}

\label{lesson:ZH-C495-zuobi}



\begin{quote}
Before the new one: say the Mandarin for an exam hall, then the Mandarin for full marks.
\end{quote}

\subsection*{You'll want to know: zuòbì}
\emph{zuòbì} — \textquotedblleft{}to cheat in an exam\textquotedblright{}.

\begin{quote}
\emph{kǎoshì de shíhou bù yīnggāi zuòbì} — You shouldn't cheat in an exam.
\end{quote}

Its characters come later; for now it is a word for the ear.

\begin{cousinweb}
\emph{zuòbì} is built from \zh{作} 'to do' and \emph{bì} 'harm', both inherited from Middle Chinese, \emph{tsak} and \emph{bjiejH} (Guangyun, in Baxter's transcription), which glosses the second 'worn out' and 'bad'.
\end{cousinweb}

\begin{grammarlens}[title={What you've built}]
One more word for saying how it was.
\end{grammarlens}

\subsection*{Your turn}
\begin{itemize}
\raggedright
  \item \emph{Say it:} \emph{zuòbì}
  \item \emph{Say it:} \emph{zuòbì}, once more
  \item \emph{Say it:} \emph{kǎoshì de shíhou bù yīnggāi zuòbì} — You shouldn't cheat in an exam.
  \item \emph{Recall:} say the Mandarin for an exam hall, then the Mandarin for full marks, then say \emph{zuòbì} again
\end{itemize}

\begin{tcolorbox}[breakable,colback=teal!4,colframe=teal!35!black,title={Before you move on}]
What does zuòbì mean? (\textquotedblleft{}To cheat in an exam\textquotedblright{}.) Say it once more.
\end{tcolorbox}
\section[\texorpdfstring{\zh{难题} (nántí)}{难题 (nántí)}]{\zh{难题} (nántí) — hard problem; difficult question}

\label{lesson:ZH-C495-nanti}



\begin{quote}
Before the new one: say the Mandarin for full marks, then the Mandarin for to cheat in an exam.
\end{quote}

\subsection*{You'll want to know: \zh{难题}}
\textbf{\zh{难题}} — \emph{nántí} — \textquotedblleft{}hard problem; difficult question\textquotedblright{}.

\begin{quote}
\textbf{\zh{这个难题我想了一个小时}} \emph{zhège nántí wǒ xiǎng le yí ge xiǎoshí} — I spent an hour thinking about this hard question.
\end{quote}

\begin{cousinweb}
\textbf{\zh{难题}} (\emph{nántí}) is built from \zh{难} 'hard' and \zh{题} 'a heading', both inherited from Middle Chinese, \emph{nan} and \emph{dej} (Guangyun, in Baxter's transcription), which glosses the first 'hard to do' and the second 'a written heading'.
\end{cousinweb}

\begin{grammarlens}[title={What you've built}]
One more word for saying how it was.
\end{grammarlens}

\subsection*{Your turn}
\begin{itemize}
\raggedright
  \item \emph{Say it:} \emph{nántí}
  \item \emph{Say it:} \emph{nántí}, once more
  \item \emph{Say it:} \emph{zhège nántí wǒ xiǎng le yí ge xiǎoshí} — I spent an hour thinking about this hard question.
  \item \emph{Recall:} say the Mandarin for full marks, then the Mandarin for to cheat in an exam, then say \emph{nántí} again
\end{itemize}

\begin{tcolorbox}[breakable,colback=teal!4,colframe=teal!35!black,title={Before you move on}]
What does \zh{难题} mean? (\textquotedblleft{}Hard problem; difficult question\textquotedblright{}.) Say it once more.
\end{tcolorbox}
